\section{Next Steps}

The dense rank-recursive theorem changes the priority order.  The immediate
task is proof verification and simplification, followed by sharpening the one
remaining approximation primitive:
\begin{enumerate}
\item independently referee Lemma~\ref{lem:retained-rank} and the telescope
  (RR.26), including all ceiling losses, bridge-conditioning denominators, and
  finite-rank cases;
\item independently verify the positive star formula and its translation from
  F\'eray's convention, then compare the backward sampler against the count
  formula on the recorded exact cases;
\item solve the tunable positive connection problem: sample the exact branch
  law (RR.9), or approximate it to arbitrary requested $\epsilon$, in time
  polynomial in $N,L,$ and $\log(1/\epsilon)$;
\item improve the constants and exponent in (RR.2), which are intentionally
  loose and become numerically useful only at very large $N$;
\item determine the broadest dense length regime for which the retained-rank
  construction works, rather than fixing $L=2N^3$;
\item restrict the local heat-bath program to parameter regimes that avoid
  Theorem~\ref{thm:balanced-bottleneck}; in particular, quantify mixing
  inside support profiles at long lengths, or establish the all-endpoint
  base bound in Theorem~\ref{thm:base-length-reduction}; for growing
  windows, implement the exact conditionals in
  Theorem~\ref{thm:growing-window-mixing};
\item compare the rank-recursive construction with implementable reversible
  gate libraries and bounded-degree expander generators before drawing
  cryptographic conclusions.
\end{enumerate}
